% !TEX root = ../Projektdokumentation.tex
\section{Entwurf} 
\label{sec:Entwurf}
In dieser Phase soll die Implementierung vorbereitet werden. Dazu wird in einem UML-Aktivitätsdiagramm 
dargestellt, wie der Synchronisationsablauf in eine Richtung aussieht. Anschließend zeigt eine Architekturübersicht, wie diese 
Systeme miteinander integagieren. Am Ende stellt ein ER-Modell dar, wie die Datenbankstruktur für die Event Persistierung aussieht. 

\subsection{UML-Aktivitätsdiagramm}
\label{sec:UML-Aktivitätsdiagramm}
Das folgende UML-Aktivitätsdiagramm für die Synchronisation in Richtung \textit{HubSpot}, zeigt
den Ablauf der Synchronisation vom \textit{user-service} ausgehend.
Angefangen mit der Eingabe der Kontaktdaten vom Nutzer wird ein Event in die Main Queue geschickt.
Anschließend konsumiert der \textit{Nightcrawler} diese Queue und sendet basierend auf den Daten innerhalb der
Event Message, einen API Upsert Aufruf an \textit{HubSpot}. Sofern der API Aufruf einen 201 Created oder 200 Ok Statuscode zurückliefert
wird das Event als erfolgreich angesehen. Andernfalls wird das Event in eine Retry Queue geschickt und
nach einer 30 sekündigen \acs{TTL} beim ersten Mal und 60 Sekunden beim zweiten Mal erneut 
ausgeführt. Wenn das zweite Mal ebenfalls nicht erfolgreich ist, wird das Event in eine \acs{DLQ} geschickt.\\
\par\smallskip
\begin{figure}[htb]
\centering
\includegraphics[scale=0.41]{\uml}\\
\end{figure}
Für genauere technische Abläufe siehe \Anhang{app:Flows}

\subsection{Architekturdesign}
\label{sec:Architekturdesign}
Die Anwendung ist als mehrschichtige Microservice-Architektur aufgebaut.
Innerhalb der Nightcrawler-Solution existiert ein eigenes Orchestrator-Projekt,
welches den Start der Anwendung, die Konfiguration sowie die Registrierung der benötigten Services übernimmt.
Dadurch werden technische Abhängigkeiten zentral verwaltet und die einzelnen Anwendungskomponenten klar voneinander getrennt.
\par\smallskip
Die fachliche Verarbeitung ist in mehrere Schichten aufgeteilt.
Eingehende Ereignisse werden zunächst von der \textit{Presentation Layer} konsumiert und anschließend an die Anwendungsschicht weitergegeben.
Dort wird der Synchronisationsablauf koordiniert.
Die \textit{Domain Layer} enthält die fachlichen Modelle und Regeln,
während die \textit{Infrastructure Layer} die technischen Anbindungen an externe Systeme wie RabbitMQ, HubSpot und PostgreSQL bereitstellt.
\par\smallskip
Abbildung~\ref{fig:ArchitekturNightcrawler} zeigt den vereinfachten Aufbau der Nightcrawler-Solution
mit Orchestrator-Projekt, Anwendungsschichten und externen Schnittstellen.
\par\smallskip
Die Anwendung wird mit Microsoft .NET 10 und C\# entwickelt.
Diese Technologie wurde gewählt, da alle Backend-Services im Unternehmen auf C\# und .NET basieren.
Dadurch kann die neue Anwendung einheitlich in die bestehende Systemlandschaft integriert und von den vorhandenen Entwicklern gewartet werden.
\par\smallskip
Ein weiterer Grund für den Einsatz von .NET ist die plattformunabhängige Laufzeitumgebung.
Die Anwendung kann dadurch sowohl lokal während der Entwicklung als auch in der Zielumgebung auf einem Linux-basierten System betrieben werden \cite.
Zusätzlich unterstützt .NET die containerbasierte Entwicklung und Bereitstellung, was zur bestehenden Infrastruktur des Unternehmens passt.
\par\smallskip
Da .NET 10 eine aktuelle LTS-Version ist, eignet sich das Framework außerdem für eine langfristig wartbare Anwendung.
Durch den großen Funktionsumfang und verfügbare NuGet-Pakete, beispielsweise für Datenbankzugriffe oder Resilienzmechanismen,
können die Anforderungen des Projektes ohne den Einsatz weiterer Programmiersprachen umgesetzt werden.


\subsection{Datenbankstruktur}
\label{sec:Datenbankstruktur}
Die Datenstruktur liegt zur Veranschaulichung im Folgenden als \acs{ERD} vor. Sie besteht aus zwei Tabellen \textit{SyncHistory} und \textit{UserContactState}, 
welche in einer 1:n Beziehung zu einander stehen. Der \textit{UserContactState} beschreibt einen Nutzer, welcher mit \textit{HubSpot} verknüpft wird und die 
\textit{SyncHistory} beschreibt alle Aktivitäten, die dieser Nutzer im Bezug auf dessen Daten gemacht hat.
\par\smallskip
Siehe \Anhang{app:ERD} für das \ac{ERD} zur \textit{nightcrawler} Datenbank, welches Attribute, Relationen und die dazugehörigen Multiplizitäten enthält. 
